% META: {"id": "1NSI-RES-EVAL-B", "chapitre": "1NSI-RESEAUX", "type_objet": "evaluation", "version": "B", "capacites": ["P-ARCH-02A", "P-ARCH-02B", "P-ARCH-02C", "P-ARCH-04A", "P-ARCH-04B"], "mode_creation": "ex_nihilo", "sources_inspiration": [], "duree_min": 45, "status": "verified"}
\section*{Évaluation B --- Réseaux}
\textit{Durée : 45 minutes. Ordinateur avec interpréteur Python autorisé.}

\baremeIndicatif{Ex. 1 : 8 pts (paquets, bit alterné) ; Ex. 2 : 8 pts (réseau, capteurs) ; Ex. 3 : 4 pts (IHM sur cahier des charges)}

\bigskip

\textbf{Exercice 1} --- \textit{Paquets et bit alterné} \hfill (8 points)

\begin{enumerate}
  \item (3 pts) Découper le message suivant en paquets de taille maximale $5$ avec
  \lstinline{decouper_en_paquets} du cours :
  \begin{center}\lstinline{"PAQUETSRESEAU"}\end{center}
  Combien de paquets sont créés, et que contiennent-ils ?
  \item (3 pts) On envoie les messages \lstinline{["M", "N", "O"]} avec le protocole du bit alterné ; la deuxième tentative globale échoue. Dérouler l'envoi et donner la liste finale des messages reçus, avec leur bit.
  \item (2 pts) Pourquoi un seul bit (0 ou 1) suffit-il à distinguer un nouveau paquet d'une retransmission, dans ce protocole ?
\end{enumerate}

% BEGIN-VERIFY
% def decouper_en_paquets(message, taille_max):
%     nb_paquets = (len(message) + taille_max - 1) // taille_max
%     paquets = []
%     for i in range(nb_paquets):
%         donnees = message[i * taille_max:(i + 1) * taille_max]
%         paquets.append({"numero": i, "total": nb_paquets, "donnees": donnees})
%     return paquets
% paquets = decouper_en_paquets("PAQUETSRESEAU", 5)
% assert paquets == [
%     {"numero": 0, "total": 3, "donnees": "PAQUE"},
%     {"numero": 1, "total": 3, "donnees": "TSRES"},
%     {"numero": 2, "total": 3, "donnees": "EAU"},
% ]
% def protocole_bit_alterne(messages, tentatives_perdues):
%     resultat = []
%     bit = 0
%     compteur = 0
%     for m in messages:
%         recu = False
%         while not recu:
%             compteur += 1
%             if compteur in tentatives_perdues:
%                 continue
%             resultat.append((bit, m))
%             recu = True
%         bit = 1 - bit
%     return resultat
% assert protocole_bit_alterne(["M", "N", "O"], {2}) == [(0, "M"), (1, "N"), (0, "O")]
% END-VERIFY

\bigskip

\textbf{Exercice 2} --- \textit{Réseau et capteurs} \hfill (8 points)

\begin{enumerate}
  \item (3 pts) On donne le réseau suivant :
  \begin{center}\lstinline|{"X": ["Y"], "Y": ["X", "Z"], "Z": ["Y"]}|\end{center}
  En utilisant \lstinline{peuvent_communiquer} du cours, vérifier que \lstinline{"X"}
  et \lstinline{"Z"} peuvent communiquer.
  \item (3 pts) Écrire une fonction \lstinline{decider_ventilation(co2, seuil=800)} qui renvoie \lstinline{True} si le taux de CO2 (en ppm) dépasse le seuil (la ventilation doit s'activer).
  \item (2 pts) Tester cette fonction pour un taux de $1000$ ppm et de $500$ ppm. Identifier, dans ce système, le capteur et l'actionneur.
\end{enumerate}

% BEGIN-VERIFY
% reseau = {"X": ["Y"], "Y": ["X", "Z"], "Z": ["Y"]}
% def peuvent_communiquer(reseau, depart, arrivee, visites=None):
%     if visites is None:
%         visites = set()
%     if depart == arrivee:
%         return True
%     visites.add(depart)
%     for v in reseau.get(depart, []):
%         if v not in visites and peuvent_communiquer(reseau, v, arrivee, visites):
%             return True
%     return False
% assert peuvent_communiquer(reseau, "X", "Z") is True
% def decider_ventilation(co2, seuil=800):
%     return co2 > seuil
% assert decider_ventilation(1000) is True
% assert decider_ventilation(500) is False
% END-VERIFY

\bigskip

\textbf{Exercice 3} --- \textit{IHM repondant a un cahier des charges} \hfill (4 points)

\textit{Cahier des charges.} Un local est equipe d'un capteur de CO2. L'IHM recoit une
commande textuelle parmi \lstinline{"AUTO"}, \lstinline{"ON"} et \lstinline{"OFF"} :
\lstinline{"ON"} force la ventilation, \lstinline{"OFF"} l'interdit, et \lstinline{"AUTO"}
rend la decision au capteur (ventilation si le taux depasse $800$ ppm). Toute autre commande
doit être refusee par une erreur explicite.

\begin{enumerate}
  \item (3 pts) Écrire les fonctions \lstinline{traiter_commande(commande)} et
  \lstinline{ventilation_active(co2)} realisant ce cahier des charges, sur le modèle du
  thermostat du cours.
  \item (1 pt) Donner trois cas de test qui, ensemble, couvrent les trois modes, et dire ce
  que chacun vérifie.
\end{enumerate}

% BEGIN-VERIFY
% COMMANDES = {"AUTO": None, "ON": True, "OFF": False}
% etat = {"override_manuel": None}
%
% def traiter_commande(commande):
%     commande = commande.strip().upper()
%     if commande not in COMMANDES:
%         raise ValueError("commande attendue : AUTO, ON ou OFF")
%     etat["override_manuel"] = COMMANDES[commande]
%
% def ventilation_active(co2, seuil=800):
%     if etat["override_manuel"] is not None:
%         return etat["override_manuel"]
%     return co2 > seuil
%
% # Mode AUTO : la decision revient au capteur, de part et d'autre du seuil.
% traiter_commande("auto")
% assert ventilation_active(1000) is True
% assert ventilation_active(500) is False
% assert ventilation_active(800) is False        # le seuil lui-meme ne declenche pas
% # Mode ON : la ventilation est forcee malgre un taux bas.
% traiter_commande("ON")
% assert ventilation_active(500) is True
% # Mode OFF : la ventilation est interdite malgre un taux eleve.
% traiter_commande("OFF")
% assert ventilation_active(1000) is False
% # Une commande hors cahier des charges est refusee explicitement.
% try:
%     traiter_commande("MAX")
% except ValueError as e:
%     assert "AUTO" in str(e)
% else:
%     raise AssertionError("la commande invalide n'a pas ete refusee")
% END-VERIFY
